\FloatBarrier
\noindent This paper examined whether a passively oriented, spin-stabilised solar sail in near-polar low Earth orbit can exhibit a bounded interval of net orbital-energy gain under eclipse-modulated solar-radiation-pressure forcing. The results show that such a regime does exist, but only over a narrow and well-defined region of the sampled LTAN--pitch space. The gain response is therefore not a generic property of passive sailing in LEO, but a constrained transient effect that appears only under specific Sun-relative geometries.
\vspace{0.5em}
\par \noindent The results also show that eclipse is the causal enabler of this transient gain regime. In the representative gain case, the inclusion of eclipse produces positive short-horizon semi-major-axis growth, whereas removing eclipse from the same initial condition eliminates that behaviour and leads instead to clear decay. Orbit-resolved SRP work diagnostics show why: eclipse does not merely weaken SRP forcing, but selectively removes unfavourable portions of the SRP cycle. Under the gain geometry, this shifts the orbit-averaged SRP work balance positive; under neutral and decaying geometries, the same shadowing instead suppresses the positive recovery branch and leaves the net SRP contribution negative. The central physical result is therefore that eclipse timing can control the sign of the orbit-averaged SRP work input.
\vspace{0.5em}
\par \noindent At the same time, the longer-horizon behaviour shows unambiguously that this effect is transient and self-limiting rather than sustained. The representative gain trajectory initially exhibits semi-major-axis increase, but this does not correspond to uniform orbit raising. Instead, the orbit evolves through apogee growth, perigee erosion, and rising eccentricity, so that the early energy advantage progressively creates the conditions that undermine it. Drag becomes increasingly concentrated near perigee, while the favourable SRP projection weakens as the orbit distorts. As a result, the late-time orbit-averaged energy balance turns negative and the semi-major axis ultimately declines below its initial value. The transient gain regime identified here must therefore be understood as a front-loaded energy imbalance, not as evidence of persistent passive orbital uplift.

\noindent The sensitivity analysis supports the same interpretation. The gain-anchor case is not a numerically isolated point, but neither is it broadly robust. At lower altitude, the response is strongly limited by density and attitude offset; at higher altitude, the same geometry remains more tolerant over the sampled envelope, but still does not become a general passive orbit-raising solution. The mechanism is therefore physically real and interpretable, yet intrinsically bounded by geometry, atmospheric conditions, and the coupled eccentricity and drag penalties that emerge during the evolution.

\noindent Accordingly, the main contribution of this paper is not to propose sustained passive orbit raising in LEO, but to establish a narrower result: a constrained transient gain regime exists, eclipse timing is its causal enabler, and the same coupled dynamics that permit the early gain also drive its later breakdown. Future work should remain tightly focused on this mechanism, particularly on better analytical bounds for the gain window, broader mapping of the geometry dependence, and controlled studies of whether limited attitude variation --- informed by optimal steering approaches such as those described by \citet{dachwaldOptimalSolarSail2005} and \citet{mcinnes1999} --- can delay the eccentricity and drag-driven breakdown without changing the fundamentally transient character of the effect.
\FloatBarrier